\begin{frame}{Indexation : motivation}
  \begin{itemize}
    \item Par défaut, un moteur de recherche essaie d'analyser les documents.
    \item Pour chaque champ, il tente de trouver le type des données  (entier, date, IP, texte en français, en anglais, etc.)
\item Acceptable pour le texte brut, mais on souhaite généralement faire mieux: ``on 
connaît toujours mieux ses données qu'ElasticSearch'' 
  \end{itemize}

\vfill

\begin{blocgauche}
Le typage implique  une \alert{analyse} et une \alert{transformation} des données
pendant un \alert{pré-traitement}. Ce dernier
  conditionne la qualité des résultats.
\end{blocgauche}

\end{frame}


\begin{comment}
\begin{frame}{Pré-traitement: l'analyse}
  \begin{itemize}
    \item Dès qu'un document est un peu complexe, on ne peut pas le découper
      arbitrairement sans appliquer un pré-traitement réfléchi
    \item sinon :
      \begin{itemize}
        \item que devient ``pomme de terre'' ?
        \item on stocke des listes très longues pour certains mots\ldots
        \item l'utilisateur doit faire très attention (``loup'', ``Loup'',
          ``loups'' sont différents)
      \end{itemize}
  \end{itemize}
  

\vfill

\begin{blocgauche}
Il faut procéder à une \alert{analyse du contenu}, ce qui suppose de savoir
à quel type de texte on a affaire (langue par exemple).
\end{blocgauche}

\end{frame}
\end{comment}

%\begin{frame}{Pré-traitement: la transformation}

%Quelques requêtes à faire sur la base des films au préalable:

%\begin{itemize}
  %\item \texttt{Matrix}, \texttt{MATRIX} et \texttt{matrix}
  %\item \texttt{title:Matrix} et \texttt{text:Matrix}
  %\item \texttt{title:reloaded} et \texttt{text:reloaded}
%\end{itemize}

%\vfill

%\begin{block}{Que se passe-t-il?}
%Si vous avez bien suivi les explications sur le schéma vous devriez
%comprendre ce qui se passe. Réflechissons ensemble...
%\end{block}
%\end{frame}

\begin{frame}{Rôle de l'analyse}

L'analyse des textes amène à effectuer une forme de 
normalisation / unification pour être moins dépendant de la
forme du texte.

\begin{itemize}
  \item un document parle de loup même si on y trouve les formes
     "loups", "Loup", "louve", etc.
  \item un document parle de travail quelle que soit la forme
     du verbe "travailler" ou de ses variantes.
  \item Jusqu'où va-t-on? Traductions (loup = wolf = lupus) ? Synonymes
      (loup  = prédateur) ? (sujet de recherches)
\end{itemize}

\vfill

\begin{blocgauche}
On applique des \alert{transformations} pour moins dépendre de la forme 
\end{blocgauche}

\end{frame}


\begin{frame}{Impact de l'analyse}


\alert{Plus on normalise, plus on diminue la précision}.
\smallskip
Car des mots distincts sont unifiés (cote, côte, côté, etc.)

\vfill

\alert{Plus on normalise, plus on améliore le rappel}.
\smallskip
Car on met en correspondance les variantes d'un même mot, d'une même
signification (conjugaisons d'un verbe).

\vfill

  \begin{defblock}{Très important}
    La même transformation doit être appliquée aux documents \alert{et} à la
    requête. 
  \end{defblock}
  \vfill
\begin{blocgauche}
Sinon, il se crée un décalage entre ce qui a été analysé et ce qui peut
      être réellement recherché
  \end{blocgauche}
  
\vfill

\end{frame}

\begin{frame}{Les phases de l'analyse}

Important: identification de quelques méta-données (la langue), prise en compte du contexte
(quels documents pour quelle application). Puis, 

\vfill

\begin{itemize}
\item \alert{Tokenization}\,: découpage du texte en ``mots''
\item \alert{Normalisation}\,: majuscules\,? acronymes\,? apostrophes\,? accents\,?

     Exemple\,: \textit{Windows} et \textit{window}, U.S.A vs USA,
     \emph{l'étudiant}  vs \emph{les étudiants}.
     
     \item \alert{Stemming} (``racinisation''), \alert{lemmatization}
     
     Prendre la racine des mots pour éviter le biais des variations (étudier,
     étudiant, étude, etc.)
     
     \item  \alert{Stop words}, quels mots garder\,?
     
     Mots très courants peu informatifs (le, un à, de).
\end{itemize}

\begin{blocgauche}
C'est de l'art et du réglage... Dans ce qui suit\,: introduction /
sensibilisation aux problèmes.
\end{blocgauche}

\end{frame}

\subsection{Tokenisation et normalisation}

\begin{comment}
\begin{frame}{Identification de la langue}

Comment trouver la langue d'un document\,?

\begin{itemize}

\item Méta information (dans l'entête HTTP p.e.): pas fiable du tout.

\item \alert{Par le jeu de caractères}, pas assez courant!

\smallskip
\figSlideWithSize{characters}{2cm}

\pause

\structure{Respectivement}: Coréen, Japonais, Maldives, Malte, Islandais.

\item Extension: \alert{séquences de caractères fréquents}, ($n$-grams)

\item Techniques d'apprentissage (\alert{classifiers})

\end{itemize}

\end{frame}
\end{comment}

\begin{frame}{Tokenisation}

\begin{defblock}{Principe}
Séparation du texte en \alert{tokens} (``mots'')
\end{defblock}

\smallskip


Pas du tout aussi facile qu'on le dirait\,!

\begin{itemize}

\item Dans certaines langues (Chinois, Japonais), les mots  \alert{ne sont pas}
séparés par des espaces.

\item Certaines langues s'écrivent de droite à gauche, de haut en bas.

\end{itemize}


\begin{blocgauche}
Que faire (et de manière \alert{cohérente}) des acronymes, élisions, nombres, unités, URL, email, etc.
\end{blocgauche}

\end{frame}

\begin{frame}{Tokenisation}

\alert{Mots composés}\,: les séparer en \emph{tokens} ou les regrouper en un seul\,?

\begin{enumerate}
   \item Anglais\,: \emph{hostname}, \emph{host-name} et \emph{host
name}, ...
  \item Français\,: Le Mans, aujourd'hui, pomme de terre, ...
  \item Allemand\,: \emph{Levensversicherungsgesellschaftsangestellter} (employé d’une
société d’assurance vie)
\end{enumerate}

\vfill

 Que faire si l'utilisateur
cherche \emph{hostname} et qu'on a normalisé en \emph{host-name}\,?

\vfill
\begin{blocgauche}
Majuscules, ponctuation\,? Une solution simple est de normaliser (minuscules, pas de ponctuation).
\end{blocgauche}

\end{frame}

\begin{frame}{Exemple pour notre petit jeu de données}

On met en minuscules, on retire la ponctuation.

\vfill

\begin{tabular}{>{\color{structure.fg}}ll}
     $d_1$&le loup est dans la bergerie\\
     $d_2$&le loup et les trois petits cochons\\
     $d_3$&les moutons sont dans la bergerie\\
     $d_4$&spider cochon spider cochon il peut marcher au plafond\\
     $d_5$&un loup a mangé un mouton les autres loups sont restés dans la bergerie\\
     $d_6$&il y a trois moutons dans le pré et un mouton dans la gueule du loup\\
     $d_7$&le cochon est à 12 euros le kilo le mouton à  10 euros le kilo\\
     $d_8$&les trois petits loups et le grand méchant cochon\\
\end{tabular}

\vfill

On considère que l'espace est le séparateur de tokens.

\end{frame}


\subsection{Racinisation}

\begin{frame}{Stemming (racine)}
\begin{defblock}{Principe}
\alert{Confondre} toutes les formes d'un mot, ou de mots apparentés, en
une seule \alert{racine}.
\end{defblock}

\smallskip

\begin{description}

\item[\alert{Stemming Morphologique.}] Retire les pluriels, marque de genre, conjugaisons, modes.

\begin{itemize}

\item Très dépendant de la langue\,: \emph{geese} pluriel de \emph{goose}, \emph{mice} de \emph{mouse}

\item Difficile à séparer d'une analyse linguistique (``Les poules du couvent couvent'', ``la petite
    brise la glace''\,:  où est le verbe\,?)
\end{itemize}

\item[\alert{Stemming lexical}] Fondre les termes proches lexicalement\,: ``politique, politicien, police (?)'' ou 
     ``université, universel, univers (?)'' 

\end{description}

\begin{blocgauche}
\alert{Stemming phonétique.}  Correction fautes de frappes, fautes orthographes
\end{blocgauche}

\end{frame}

\begin{frame}{Exemple de stemming}

On retire les pluriels, on met le verbe à l'infinitif.

\vfill

\begin{tabular}{>{\color{structure.fg}}ll}
     $d_1$&le loup etre dans la bergerie\\
     $d_2$&le loup et les trois petit cochon\\
     $d_3$&les moutons etre dans la bergerie\\
     $d_4$&spider cochon spider cochon il pouvoir marcher au plafond\\
     $d_5$&un loup avoir manger un mouton les autres loups etre rester dans la bergerie\\
     $d_6$&il y avoir trois mouton dans le pre et un mouton dans la gueule du loup\\
     $d_7$&le cochon etre a 12 euro le kilo le mouton a 10 euro le kilo\\
     $d_8$&les trois petit loup et le grand mechant cochon\\
\end{tabular}

\end{frame}


\subsection{Mots vides, synonymes, cas particuliers...}


\begin{frame}{Suppression des \textit{Stop Words}}

\begin{defblock}{Principe}
On retire les mots porteurs d'une information faible afin de limiter le stockage.
\end{defblock}

\begin{description}
\item[articles:] \emph{le}, \emph{le}, \emph{ce}, etc.

\item[verbes ``fonctionnels''] \emph{être}, \emph{avoir}, \emph{faire}, etc.

\item[conjunctions:] \emph{that}, \emph{and}, etc.

\item[etc.]
\end{description}


\begin{blocgauche}
\remark{
Maintenant moins utilisé car (i) espace de stockage peu coûteux et (ii) pose
d'autres problèmes (``pomme de terre'', ``Let it be'', ``Stade de France'')
}
\end{blocgauche}

\end{frame}



\begin{frame}{Autres problèmes, en vrac}


\alert{Majuscules / minuscules}

Lyonnaise des Eaux, Société Générale, etc.

\vfill
 \alert{Acronymes}
 
   CAT = \textit{cat} ou \textit{Caterpillar Inc.}\,?  M.A.A.F ou MAAF ou Mutuelle ... \,?

\vfill

\alert{Dates, chiffres}

   Monday 24, August, 1572 -- 24/08/1572 -- 24 août 1572
    10000 ou 10,000.00 ou 10,000.00
     
     
     \vfill
     
     \alert{Accents, ponctuation}
     
     résumé ou résume ou resume... 
     
\begin{blocgauche}
 \remark{Dans tous les cas, les même règles de transformation s'appliquent aux documents ET à la requête. }
\end{blocgauche}

 \end{frame}


\begin{frame}{Exemple avec suppression des \textit{stop words}}

Voici une solution possible.

\vfill

\begin{tabular}{>{\color{structure.fg}}ll}
     $d_1$&loup etre bergerie\\
     $d_2$&loup  trois petit cochon\\
     $d_3$&mouton etre bergerie\\
     $d_4$&spider cochon spider cochon pouvoir marcher plafond\\
     $d_5$&loup avoir manger mouton  autres loups etre rester bergerie\\
     $d_6$&avoir trois mouton pre mouton gueule loup\\
     $d_7$&cochon etre 12 euro kilo mouton 10 euro kilo\\
     $d_8$&trois petit loup grand mechant cochon\\
\end{tabular}

\vfill

On a gardé les verbes fonctionnels (être, avoir).
\end{frame}



\begin{frame}{Résumé}

Ce qui précède est simplement une sensibilisation. Retenir:

\begin{itemize}
 \item Installer un moteur de recherche, \alert{c'est facile}
 \item Obtenir des résultats satisfaisants, \alert{c'est beaucoup plus dur}.
 \end{itemize}

\vfill

Les questions à se poser.

\begin{itemize}
 \item Quelle est la langue de mes documents?
 \item Quels types de recherche vont être soumis par mes utilisateurs?
 \item Comme évaluer et améliorer mes résultats?
 \end{itemize}


\vfill
\begin{blocgauche}
\alert{Et on peut aller plus loin}: facettes, synonymes, traductions, retours
utilisateurs, etc.
\end{blocgauche}

\end{frame}


\end{document}

\subsection{Analyse avec ElasticSearch}

\begin{frame}[fragile]{Définir une chaîne d'analyse avec Solr}

\begin{itemize}
	\item Les analyseurs sont composés d'un tokenizer, et TokenFilters (0 ou plus)
	\item Le tokenizer peut être précédé de CharFilters
	 \item les filtres (filter) examinent les tokens un par un et décident de les conserver, de les remplacer par un ou plusieurs autres;
	 \item l’analyseur est une chaîne de traitement (pipeline) constituée de tokeniseurs et de filtres.
\end{itemize}

\end{frame}

\begin{frame}[fragile]{Définir une chaîne d'analyse avec Solr}
\begin{minted}[frame=lines, baselinestretch=.9, fontsize=\tiny, framesep=2mm]{json}
{
  "settings": {
    "analysis": {
      "filter": {
        "custom_english_stemmer": {
          "type": "stemmer",
          "name": "english"
        }
      },
      "analyzer": {
        "custom_lowercase_stemmed": {
          "tokenizer": "standard",
          "filter": [
            "lowercase",
            "custom_english_stemmer"
          ]
        }
      }
    }
  },
  "mappings": {
    "test": {
      "properties": {
        "text": {
          "type": "string",
          "analyzer": "custom_lowercase_stemmed"
        }
      }
    }
  }
}'
\end{minted}
	
\end{frame}



\begin{frame}[fragile]{tester les analyseurs}
\begin{minted}[frame=lines, baselinestretch=.9, fontsize=\tiny, framesep=2mm]{json}
  {                                                                               
    "analyzer": "english",                                                        
    "text": "j'aime les couleurs de l'arbre durant l'été"                         
  }                                                                               
  {                                                                               
    "analyzer": "french",                                                         
    "text": "j'aime les couleurs de l'arbre durant l'été"                         
  }                                                                               
  {                                                                               
    "analyzer": "standard",                                                       
    "filter":  [ "lowercase", "asciifolding" ],                                   
    "text": "j'aime les COULEURS de l'arbre durant l'été"                         
  }   
\end{minted}

\end{frame}

